\documentclass{ximera}
%verbeterd 16/11/26

\title{Vragen potentiële energie, potentiaal en spanning}
\author{Daan Lipkens}

\begin{document}
\begin{abstract}
\end{abstract}
\maketitle
\label{vragen:potentiaal}
\section{Vragen potentiële energie, potentiaal en spanning}

%———————————————————————————————————————
% VRAAG 1: TEKEN VAN Epot BIJ VERPLAATSING
\begin{question}
    Een negatieve lading beweegt in een homogeen elektrisch veld van de negatieve naar de positieve plaat. Wat gebeurt er met haar potentiële elektrische energie?

    \begin{multipleChoice}
        \choice{Ze stijgt.}
        \choice[correct]{Ze daalt.}
        \choice{Ze blijft gelijk.}
    \end{multipleChoice}

    \begin{solution}\nl
        Een negatieve lading wordt aangetrokken door de positieve plaat, dus de elektrische kracht werkt in de zin van de verplaatsing. De beweging is spontaan: het veld verricht positieve arbeid. Omdat $\Delta E_{\text{pot}} = -W$, daalt haar potentiële energie.
    \end{solution}
\end{question}
%———————————————————————————————————————

%———————————————————————————————————————
% VRAAG 2: RICO E_pot(d)-GRAFIEK
\begin{question}
    De grafiek van $E_{\text{pot}}$ in functie van $d$ voor een lading $Q$ in een homogeen veld is een \textbf{dalende} rechte door de oorsprong. Wat kan je hieruit besluiten over het teken van $Q$?

    \begin{multipleChoice}
        \choice[correct]{$Q$ is negatief.}
        \choice{$Q$ is positief.}
        \choice{Er zijn te weinig gegevens om dit te bepalen.}
    \end{multipleChoice}

    \begin{hint}
        De richtingscoëfficiënt van de rechte $E_{\text{pot}}(d)$ is gelijk aan $Q \cdot E$. Het elektrisch veld $E$ is per definitie altijd positief.
    \end{hint}

    \begin{solution}\nl
        De richtingscoëfficiënt van de rechte $E_{\text{pot}}(d)$ is $Q \cdot E$. Omdat $E$ steeds positief is, kan een dalende (negatieve) richtingscoëfficiënt enkel ontstaan als $Q$ negatief is.
    \end{solution}
\end{question}
%———————————————————————————————————————

%———————————————————————————————————————
% VRAAG 3: VERGELIJKEN VAN TWEE PUNTEN OP EEN VELDLIJN
\begin{question}
    Twee punten $1$ en $2$ liggen op eenzelfde veldlijn van een homogeen elektrisch veld, waarbij punt $2$ dichter bij de positieve plaat ligt dan punt $1$. Vergelijk de potentiaal in beide punten.

    \begin{multipleChoice}
        \choice{$V(1) > V(2)$}
        \choice[correct]{$V(1) < V(2)$}
        \choice{$V(1) = V(2)$}
    \end{multipleChoice}

    \begin{solution}\nl
        De potentiaal in een homogeen veld is $V = E \cdot d$, met $d$ de afstand tot de negatieve plaat. Omdat punt $2$ dichter bij de positieve plaat ligt, is $d_2 > d_1$ en dus ook $V(2) > V(1)$.
    \end{solution}
\end{question}
%———————————————————————————————————————

%———————————————————————————————————————
% VRAAG 6: ENERGIE-EENHEID VERBAND MET SPANNING
\begin{question}
    Toon aan, vertrekkend van de definities van spanning en arbeid, dat de eenheid volt (V) gelijk is aan joule per coulomb (J/C).

    \begin{solution}\nl
        Uit $W = Q \cdot U$ volgt dat $U = W/Q$. De eenheid van arbeid $W$ is joule (J) en de eenheid van lading $Q$ is coulomb (C). Bijgevolg is de eenheid van spanning $\mathrm{J/C}$. Die eenheid noemen we volt: $1 \, \mathrm{V} = 1 \, \mathrm{J/C}$. Voor de potentiaal schrijven we de eenheid als $\mathrm{J/C}$, voor de spanning als $\mathrm{V}$.
    \end{solution}
\end{question}
%———————————————————————————————————————

%———————————————————————————————————————
% VRAAG 9: DIPOOL, ELEKTRON LOSLATEN
\begin{question}
    De tekening toont de veldlijnen tussen een positieve en een negatieve lading. Punt $A$ ligt op de middelste veldlijn. Een elektron wordt in $A$ in rust losgelaten.

    \begin{center}
        \begin{tikzpicture}
            \poslading{P}{-1.5,0}
            \neglading{N}{1.5,0}
            \foreach \a in {0,40,-40,75,-75} {
                \draw[Eveldlijn,fill=none] ($(P)+(\a:0.3)$) to[out=\a,in={180-\a},looseness=1.2] ($(N)+({180-\a}:0.3)$);
            }
            % veldlijnen aan de buitenkant (berekend in Python)
                \draw[Eveldlijn,fill=none] plot[smooth] coordinates {(-1.80,0.00) (-2.21,0.00) (-2.61,0.00) (-3.01,0.00) (-3.41,0.00)};
                \draw[Eveldlijn,fill=none] plot[smooth] coordinates {(-1.75,0.17) (-2.08,0.41) (-2.40,0.66) (-2.70,0.91) (-3.01,1.17) (-3.30,1.45) (-3.40,1.55)};
                \draw[Eveldlijn,fill=none] plot[smooth] coordinates {(-1.75,-0.17) (-2.08,-0.41) (-2.40,-0.66) (-2.70,-0.91) (-3.01,-1.17) (-3.30,-1.45) (-3.40,-1.55)};
                \draw[Eveldlijn,fill=none] plot[smooth] coordinates {(3.41,0.00) (3.01,0.00) (2.61,0.00) (2.21,0.00) (1.80,0.00)};
                \draw[Eveldlijn,fill=none] plot[smooth] coordinates {(3.40,1.55) (3.30,1.45) (3.01,1.17) (2.70,0.91) (2.40,0.66) (2.08,0.41) (1.75,0.17)};
                \draw[Eveldlijn,fill=none] plot[smooth] coordinates {(3.40,-1.55) (3.30,-1.45) (3.01,-1.17) (2.70,-0.91) (2.40,-0.66) (2.08,-0.41) (1.75,-0.17)};
            \fill (-0.5,0) circle (1.5pt) node[above] {$A$};
        \end{tikzpicture}
    \end{center}

    Wat gebeurt er met het elektron?

    \begin{multipleChoice}
        \choice{Het beweegt naar rechts, naar de negatieve lading. Zijn potentiële energie daalt.}
        \choice[correct]{Het beweegt naar links, naar de positieve lading. Zijn potentiële energie daalt.}
        \choice{Het beweegt naar links, naar de positieve lading. Zijn potentiële energie stijgt.}
        \choice{Het blijft in rust.}
    \end{multipleChoice}

    \begin{hint}
        In welke zin wijst het elektrisch veld in $A$? En in welke zin werkt de kracht op een \emph{negatieve} lading ten opzichte van het veld?
    \end{hint}

    \begin{solution}\nl
        Het veld in $A$ wijst naar rechts (van $+$ naar $-$). Op een negatieve lading werkt de kracht in de tegengestelde zin: naar links. Het elektron wordt dus aangetrokken door de positieve lading en beweegt spontaan in die zin. Het veld verricht positieve arbeid, dus zijn potentiële energie daalt.
    \end{solution}
\end{question}
%———————————————————————————————————————

%———————————————————————————————————————
% VRAAG 10: HOMOGEEN VELD, TWEE PUNTEN OP DEZELFDE AFSTAND VAN DE PLAAT
\begin{question}
    In een homogeen elektrisch veld tussen twee platen liggen de punten $A$ en $B$ op dezelfde afstand van de negatieve plaat, maar op een verschillende hoogte.

    \begin{center}
        \begin{tikzpicture}
            \draw[thick] (0,-1.5) -- (0,1.5);
            \draw[thick] (4,-1.5) -- (4,1.5);
            \foreach \y in {-1.2,-0.6,0,0.6,1.2} {
                \node[left] at (-0.1,\y) {$-$};
                \node[right] at (4.1,\y) {$+$};
                \draw[Eveldlijn] (4,\y) -- (0,\y);
            }
            \fill (2,0.9) circle (1.5pt) node[above right=1pt] {$A$};
            \fill (2,-0.8) circle (1.5pt) node[below right=1pt] {$B$};
        \end{tikzpicture}
    \end{center}

    Wat is de spanning $U_{AB}$ tussen deze twee punten? Verklaar.

    \begin{solution}\nl
        De spanning is $0 \, \mathrm{V}$. In een homogeen veld geldt $V = E \cdot d$, met $d$ de afstand tot de negatieve plaat. Omdat $A$ en $B$ even ver van de negatieve plaat liggen, is $V(A) = V(B)$ en dus $U_{AB} = V(A) - V(B) = 0 \, \mathrm{V}$. Verplaats je een lading van $A$ naar $B$, dan verricht het veld geen arbeid.
    \end{solution}
\end{question}
%———————————————————————————————————————

\end{document}